% Options for packages loaded elsewhere
\PassOptionsToPackage{unicode}{hyperref}
\PassOptionsToPackage{hyphens}{url}
\documentclass[
]{article}
\usepackage{xcolor}
\usepackage{amsmath,amssymb}
\setcounter{secnumdepth}{-\maxdimen} % remove section numbering
\usepackage{iftex}
\ifPDFTeX
  \usepackage[T1]{fontenc}
  \usepackage[utf8]{inputenc}
  \usepackage{textcomp} % provide euro and other symbols
\else % if luatex or xetex
  \usepackage{unicode-math} % this also loads fontspec
  \defaultfontfeatures{Scale=MatchLowercase}
  \defaultfontfeatures[\rmfamily]{Ligatures=TeX,Scale=1}
\fi
\usepackage{lmodern}
\ifPDFTeX\else
  % xetex/luatex font selection
\fi
% Use upquote if available, for straight quotes in verbatim environments
\IfFileExists{upquote.sty}{\usepackage{upquote}}{}
\IfFileExists{microtype.sty}{% use microtype if available
  \usepackage[]{microtype}
  \UseMicrotypeSet[protrusion]{basicmath} % disable protrusion for tt fonts
}{}
\makeatletter
\@ifundefined{KOMAClassName}{% if non-KOMA class
  \IfFileExists{parskip.sty}{%
    \usepackage{parskip}
  }{% else
    \setlength{\parindent}{0pt}
    \setlength{\parskip}{6pt plus 2pt minus 1pt}}
}{% if KOMA class
  \KOMAoptions{parskip=half}}
\makeatother
\setlength{\emergencystretch}{3em} % prevent overfull lines
\providecommand{\tightlist}{%
  \setlength{\itemsep}{0pt}\setlength{\parskip}{0pt}}
\usepackage{bookmark}
\IfFileExists{xurl.sty}{\usepackage{xurl}}{} % add URL line breaks if available
\urlstyle{same}
\hypersetup{
  hidelinks,
  pdfcreator={LaTeX via pandoc}}

\author{}
\date{}

\begin{document}

\section{Research Methodology: Methods and Techniques (Exam-Oriented
Notes)}\label{research-methodology-methods-and-techniques-exam-oriented-notes}

\emph{Based on the Book by C.R. Kothari (Second Revised Edition)}
{[}2{]}

\begin{center}\rule{0.5\linewidth}{0.5pt}\end{center}

\subsection{Chapter 1: Research Methodology: An
Introduction}\label{chapter-1-research-methodology-an-introduction}

\subsubsection{Meaning of Research}\label{meaning-of-research}

\begin{itemize}
\tightlist
\item
  \textbf{Definition}: A scientific and systematic search for pertinent
  information on a specific topic {[}20{]}.
\item
  \textbf{Concept}: It is an art of scientific investigation, a
  ``systematized effort to gain new knowledge'' {[}20{]}, and an
  original contribution to the existing stock of knowledge making for
  its advancement {[}21{]}.
\item
  \textbf{Essence}: A movement from the known to the unknown; a voyage
  of discovery {[}20{]}.
\end{itemize}

\subsubsection{Objectives of Research}\label{objectives-of-research}

\begin{itemize}
\tightlist
\item
  \textbf{Core Purpose}: To discover answers to questions through the
  application of scientific procedures; to find out the truth which is
  hidden and has not been discovered as yet {[}23{]}.
\item
  \textbf{Classifications}:

  \begin{enumerate}
  \def\labelenumi{\arabic{enumi}.}
  \tightlist
  \item
    \emph{Exploratory or Formulative}: Gain familiarity with a
    phenomenon or achieve new insights into it {[}23{]}.
  \item
    \emph{Descriptive}: Portray accurately the characteristics of a
    particular individual, situation, or group {[}24{]}.
  \item
    \emph{Diagnostic}: Determine the frequency with which something
    occurs or with which it is associated with something else {[}24{]}.
  \item
    \emph{Hypothesis-Testing}: Test a hypothesis of a causal
    relationship between variables {[}24{]}.
  \end{enumerate}
\end{itemize}

\subsubsection{Motivation in Research}\label{motivation-in-research}

\begin{itemize}
\tightlist
\item
  \textbf{Key Drivers}:

  \begin{itemize}
  \tightlist
  \item
    Desire to get a research degree along with its consequential
    benefits {[}24{]}.
  \item
    Desire to face the challenge in solving the unsolved problems
    {[}24{]}.
  \item
    Desire to get intellectual joy of doing some creative work {[}24{]}.
  \item
    Desire to be of service to society {[}24{]}.
  \item
    Desire to get respectability {[}24{]}.
  \end{itemize}
\end{itemize}

\subsubsection{Types of Research}\label{types-of-research}

\begin{itemize}
\tightlist
\item
  \textbf{Descriptive vs.~Analytical}:

  \begin{itemize}
  \tightlist
  \item
    \emph{Descriptive}: Surveys and fact-finding enquiries of different
    kinds. The researcher has no control over variables; they can only
    report what has happened or is happening (Ex-post facto research)
    {[}25{]}.
  \item
    \emph{Analytical}: The researcher has to use facts or information
    already available, and analyze these to make a critical evaluation
    {[}25{]}.
  \end{itemize}
\item
  \textbf{Applied vs.~Fundamental}:

  \begin{itemize}
  \tightlist
  \item
    \emph{Applied}: Aims at finding a solution for an immediate problem
    facing a society or an industrial/business organisation {[}25{]}.
  \item
    \emph{Fundamental (Basic/Pure)}: Mainly concerned with
    generalizations and with the formulation of a theory (``gathering
    knowledge for knowledge's sake'') {[}25{]}.
  \end{itemize}
\item
  \textbf{Quantitative vs.~Qualitative}:

  \begin{itemize}
  \tightlist
  \item
    \emph{Quantitative}: Based on the measurement of quantity or amount;
    applicable to phenomena that can be expressed in terms of quantity
    {[}26{]}.
  \item
    \emph{Qualitative}: Concerned with qualitative phenomenon, i.e.,
    involving quality or kind (e.g., Motivation Research to investigate
    reasons for human behavior using projective techniques) {[}26{]}.
  \end{itemize}
\item
  \textbf{Conceptual vs.~Empirical}:

  \begin{itemize}
  \tightlist
  \item
    \emph{Conceptual}: Related to some abstract idea(s) or theory;
    generally used by philosophers and thinkers to develop new concepts
    or to reinterpret existing ones {[}27{]}.
  \item
    \emph{Empirical (Experimental)}: Relies on experience or observation
    alone, often without due regard for system and theory. It is
    database-driven, coming up with conclusions which are capable of
    being verified by observation or experiment {[}27{]}.
  \end{itemize}
\item
  \textbf{Other Types}:

  \begin{itemize}
  \tightlist
  \item
    \emph{One-time vs.~Longitudinal}: Confined to a single time-period
    vs.~carried on over several time-periods {[}27{]}.
  \item
    \emph{Field-setting vs.~Laboratory vs.~Simulation}: Based on the
    environment in which it is carried out {[}27{]}.
  \item
    \emph{Clinical vs.~Diagnostic}: Case-study or indepth approaches to
    reach basic causal relations {[}27{]}.
  \item
    \emph{Exploratory vs.~Formalized}: Development of hypotheses
    vs.~substantial structure with specific hypotheses to be tested
    {[}27{]}.
  \item
    \emph{Historical}: Utilizes historical sources like documents,
    remains, etc. to study events/ideas of the past {[}27{]}.
  \item
    \emph{Conclusion-oriented vs.~Decision-oriented}: Free to
    conceptualize/redesign as one proceeds vs.~geared to the needs of a
    decision-maker {[}27{]}.
  \end{itemize}
\end{itemize}

\subsubsection{Research Approaches}\label{research-approaches}

\begin{itemize}
\tightlist
\item
  \textbf{Two Basic Approaches}:

  \begin{enumerate}
  \def\labelenumi{\arabic{enumi}.}
  \tightlist
  \item
    \emph{Quantitative}: Generation of data in quantitative form which
    can be subjected to rigorous quantitative analysis {[}28{]}.
    Sub-classified into:

    \begin{itemize}
    \tightlist
    \item
      \emph{Inferential}: Survey-based approach where a sample is
      studied to infer population characteristics {[}28{]}.
    \item
      \emph{Experimental}: Greater control over the environment; some
      variables are manipulated to observe their effect {[}28{]}.
    \item
      \emph{Simulation}: Construction of an artificial environment
      within which relevant information/data can be generated {[}28{]}.
    \end{itemize}
  \item
    \emph{Qualitative}: Subjective assessment of attitudes, opinions,
    and behavior. Uses techniques like focus group interviews,
    projective techniques, and depth interviews {[}26, 28{]}.
  \end{enumerate}
\end{itemize}

\subsubsection{Significance of Research}\label{significance-of-research}

\begin{itemize}
\tightlist
\item
  \textbf{Progress \& Inquiry}: ``All progress is born of inquiry. Doubt
  is often better than overconfidence, for it leads to inquiry, and
  inquiry leads to invention'' (Hudson Maxim) {[}29{]}.
\item
  \textbf{Inductive Thinking}: Inculcates scientific and inductive
  thinking and promotes logical habits of thinking {[}29{]}.
\item
  \textbf{Government Policy}: Serves as a vital tool for economic and
  social policy. Government research operates in three phases: (i)
  Investigation of economic structure, (ii) Diagnosis of events taking
  place, and (iii) Prognosis (prediction of future developments)
  {[}30{]}.
\item
  \textbf{Operational Problems}: Provides a quantitative basis for
  decisions in business and industry {[}29, 30{]}.
\item
  \textbf{Social Value}: Serves as a career path (degree), source of
  livelihood (professionals), outlet for new ideas (philosophers),
  developer of style (literary men), and generalizer of new theories
  (intellectuals) {[}31, 32{]}.
\end{itemize}

\subsubsection{Research Methods versus
Methodology}\label{research-methods-versus-methodology}

\begin{itemize}
\tightlist
\item
  \textbf{Research Methods}: The methods/techniques used by the
  researcher to perform research operations (grouped as: data
  collection, statistical tools for establishing relationships, and
  methods to evaluate accuracy) {[}34, 35{]}.
\item
  \textbf{Research Methodology}: A way to systematically solve the
  research problem. It is the science of studying how research is done
  scientifically {[}36{]}. It includes not only the methods but also the
  logic behind them, their assumptions, and the criteria for choosing
  specific techniques {[}36{]}.
\end{itemize}

\subsubsection{Research and Scientific
Method}\label{research-and-scientific-method}

\begin{itemize}
\tightlist
\item
  \textbf{Philosophy}: Scientific method is the pursuit of truth as
  determined by logical considerations {[}37{]}.
\item
  \textbf{Unity of Science}: ``The unity of all sciences consists alone
  in its methods, not its material'' (Karl Pearson) {[}37{]}.
\item
  \textbf{Aims}: To achieve a systematic interrelation of facts by
  experimentation, observation, and logical arguments from accepted
  postulates {[}37{]}.
\end{itemize}

\subsubsection{Importance of Knowing How Research is
Done}\label{importance-of-knowing-how-research-is-done}

\begin{itemize}
\tightlist
\item
  \textbf{Training}: Provides training in gathering material,
  card-indexing, fieldwork, statistics, questionnaires, and
  interpretation {[}40{]}.
\item
  \textbf{Objectivity}: Helps develop disciplined thinking or a `bent of
  mind' to observe fields objectively {[}41{]}.
\item
  \textbf{Intelligent Consumption}: Enables the user to intelligently
  evaluate and judge the adequacy of research results and make rational
  decisions {[}42{]}.
\end{itemize}

\subsubsection{Research Process}\label{research-process}

\begin{itemize}
\tightlist
\item
  \textbf{Flow of Activities}:

  \begin{enumerate}
  \def\labelenumi{\arabic{enumi}.}
  \tightlist
  \item
    \emph{Formulating the research problem} {[}44{]}.
  \item
    \emph{Extensive literature survey} {[}44{]}.
  \item
    \emph{Development of working hypotheses} {[}44{]}.
  \item
    \emph{Preparing the research design} {[}44{]}.
  \item
    \emph{Determining sample design} {[}44{]}.
  \item
    \emph{Collecting the data} {[}44{]}.
  \item
    \emph{Execution of the project} {[}44{]}.
  \item
    \emph{Analysis of data} {[}44{]}.
  \item
    \emph{Hypothesis testing} {[}44{]}.
  \item
    \emph{Generalisations and interpretation} {[}44{]}.
  \item
    \emph{Preparation of the report or presentation of the results}
    {[}44{]}.
  \end{enumerate}
\item
  \emph{Note}: The steps overlap continuously rather than following a
  strict sequence {[}44{]}.
\end{itemize}

\subsubsection{Criteria of Good
Research}\label{criteria-of-good-research}

\begin{itemize}
\tightlist
\item
  \textbf{Scientific Standards}:

  \begin{enumerate}
  \def\labelenumi{\arabic{enumi}.}
  \tightlist
  \item
    Purpose should be clearly defined and common concepts used {[}63{]}.
  \item
    Research procedure should be described in sufficient detail to
    permit replication {[}63{]}.
  \item
    Procedural design should be carefully planned to yield objective
    results {[}63{]}.
  \item
    Flaws in design should be reported frankly along with estimates of
    their effects {[}63{]}.
  \item
    Analysis of data should be adequate and methods of analysis
    appropriate; validity/reliability checked {[}64{]}.
  \item
    Conclusions should be confined to those justified by data {[}64{]}.
  \item
    Integrity and experience of the researcher warrant greater
    confidence {[}64{]}.
  \end{enumerate}
\item
  \textbf{Qualities}:

  \begin{itemize}
  \tightlist
  \item
    \emph{Systematic}: Structured with specified steps {[}65{]}.
  \item
    \emph{Logical}: Guided by logical reasoning (induction/deduction)
    {[}65{]}.
  \item
    \emph{Empirical}: Related to concrete data from real situations
    {[}66{]}.
  \item
    \emph{Replicable}: Allows results to be verified by repeating the
    study {[}66{]}.
  \end{itemize}
\end{itemize}

\subsubsection{Problems Encountered by Researchers in
India}\label{problems-encountered-by-researchers-in-india}

\begin{itemize}
\tightlist
\item
  \textbf{Major Bottlenecks}:

  \begin{enumerate}
  \def\labelenumi{\arabic{enumi}.}
  \tightlist
  \item
    Lack of scientific training in research methodology {[}67{]}.
  \item
    Insufficient interaction between university departments and business
    units {[}67{]}.
  \item
    Sacrosanct concept of secrecy in business, making data collection
    difficult {[}67{]}.
  \item
    Overlap and duplication of studies due to lack of information
    {[}68{]}.
  \item
    Absence of a unified code of conduct for researchers {[}68{]}.
  \item
    Difficulty of obtaining timely secretarial and computerial
    assistance {[}69{]}.
  \item
    Unsatisfactory library management and delays in tracing resources
    {[}69{]}.
  \item
    Delays in receiving governmental publications and raw statistics
    {[}70{]}.
  \item
    Ambiguities in conceptualization and definitions {[}71{]}.
  \end{enumerate}
\end{itemize}

\begin{center}\rule{0.5\linewidth}{0.5pt}\end{center}

\subsection{Chapter 2: Defining the Research
Problem}\label{chapter-2-defining-the-research-problem}

\subsubsection{What is a Research
Problem?}\label{what-is-a-research-problem}

\begin{itemize}
\tightlist
\item
  \textbf{Definition}: A difficulty experienced by a researcher
  (individual/group) in a practical or theoretical situation, seeking a
  solution {[}74{]}.
\item
  \textbf{Components of a Problem}:

  \begin{enumerate}
  \def\labelenumi{\arabic{enumi}.}
  \tightlist
  \item
    An individual or a group experiencing some difficulty {[}77{]}.
  \item
    Objective(s) to be attained {[}77{]}.
  \item
    Alternative means (courses of action) for obtaining the objective(s)
    {[}77{]}.
  \item
    Doubt in the mind of the researcher regarding the relative
    efficiency of alternatives {[}78{]}.
  \item
    An environment to which the difficulty pertains {[}78{]}.
  \end{enumerate}
\end{itemize}

\subsubsection{Selecting the Problem}\label{selecting-the-problem}

\begin{itemize}
\tightlist
\item
  \textbf{Guiding Principles}:

  \begin{itemize}
  \tightlist
  \item
    Avoid overdone subjects where throwing new light is difficult
    {[}80{]}.
  \item
    Controversial subjects should not be chosen by an average researcher
    {[}80{]}.
  \item
    Avoid too narrow or too vague problems {[}80{]}.
  \item
    The subject must be familiar, feasible, and within reach of
    resources {[}80{]}.
  \item
    Consider researcher's background, training, budget, time, and
    participant cooperation {[}81{]}.
  \item
    Precede selection with a preliminary feasibility study {[}82{]}.
  \end{itemize}
\end{itemize}

\subsubsection{Necessity of Defining the
Problem}\label{necessity-of-defining-the-problem}

\begin{itemize}
\tightlist
\item
  \textbf{Value}: ``A problem clearly stated is a problem half solved''
  {[}83{]}.
\item
  \textbf{Function}: Limits scope, separates relevant from irrelevant
  data, and dictates the research design and analytical techniques to be
  used {[}83, 84{]}.
\end{itemize}

\subsubsection{Technique Involved in Defining a
Problem}\label{technique-involved-in-defining-a-problem}

\begin{itemize}
\tightlist
\item
  \textbf{Sequential Steps}:

  \begin{enumerate}
  \def\labelenumi{\arabic{enumi}.}
  \tightlist
  \item
    \emph{Statement of the problem in a general way}: Expressing the
    practical/intellectual interest with pilot survey if needed
    {[}86{]}.
  \item
    \emph{Understanding the nature of the problem}: Discussing its
    origin and objectives with those who raised it {[}87{]}.
  \item
    \emph{Surveying the available literature}: Reviewing conceptual and
    empirical studies to find theoretical gaps and narrow the focus
    {[}88{]}.
  \item
    \emph{Developing the ideas through discussions (experience survey)}:
    Seeking advice from experienced colleagues to sharpen focus
    {[}89{]}.
  \item
    \emph{Rephrasing the research problem}: Defining technical terms,
    making basic assumptions explicit, and putting the problem in highly
    specific operational terms {[}90, 92{]}.
  \end{enumerate}
\end{itemize}

\subsubsection{An Illustration}\label{an-illustration}

\begin{itemize}
\tightlist
\item
  \emph{General}: ``Why is productivity in Japan so much higher than in
  India?'' {[}93{]}
\item
  \emph{Narrowed}: ``What factors were responsible for the higher labour
  productivity of Japan's manufacturing industries during the decade
  1971 to 1980 relative to India's manufacturing industries?'' {[}94{]}
\item
  \emph{Operational}: ``To what extent did labour productivity in 1971
  to 1980 in Japan exceed that of India in respect of 15 selected
  manufacturing industries? What factors were responsible for the
  productivity differentials between the two countries by industries?''
  {[}94{]}
\end{itemize}

\subsubsection{Conclusion}\label{conclusion}

\begin{itemize}
\tightlist
\item
  Defining a problem follows a sequential pattern of clarifying
  ambiguities, resulting in a realistic, analytically meaningful, and
  operationally viable proposition {[}96{]}.
\end{itemize}

\begin{center}\rule{0.5\linewidth}{0.5pt}\end{center}

\subsection{Chapter 3: Research Design}\label{chapter-3-research-design}

\subsubsection{Meaning of Research
Design}\label{meaning-of-research-design}

\begin{itemize}
\tightlist
\item
  \textbf{Definition}: The blueprint for the collection, measurement,
  and analysis of data; arrangement of conditions to combine relevance
  to research purpose with economy in procedure {[}99{]}.
\item
  \textbf{Parts of Research Design}:

  \begin{itemize}
  \tightlist
  \item
    \emph{Sampling design}: Selection of items to be observed {[}101{]}.
  \item
    \emph{Observational design}: Conditions under which observations are
    made {[}102{]}.
  \item
    \emph{Statistical design}: Quantitative choices regarding sample
    size and analysis {[}102{]}.
  \item
    \emph{Operational design}: Techniques to carry out specified
    procedures {[}102{]}.
  \end{itemize}
\end{itemize}

\subsubsection{Need for Research Design}\label{need-for-research-design}

\begin{itemize}
\tightlist
\item
  \textbf{Function}: Facilitates smooth sailing of research operations,
  maximizing information while minimizing expenditures of time, effort,
  and money {[}104{]}.
\end{itemize}

\subsubsection{Features of a Good
Design}\label{features-of-a-good-design}

\begin{itemize}
\tightlist
\item
  \textbf{Attributes}: Flexible, appropriate, efficient, and economical
  {[}105{]}.
\item
  \textbf{Standard}: Minimizes bias, maximizes reliability, and yields
  the smallest experimental error {[}105{]}.
\end{itemize}

\subsubsection{Important Concepts Relating to Research
Design}\label{important-concepts-relating-to-research-design}

\begin{itemize}
\tightlist
\item
  \textbf{Dependent \& Independent Variables}:

  \begin{itemize}
  \tightlist
  \item
    \emph{Variable}: A concept which can take on different quantitative
    values {[}107{]}.
  \item
    \emph{Continuous Variable}: Can assume any numerical value within a
    range (e.g., age) {[}108{]}.
  \item
    \emph{Discrete Variable}: Expressed only in integers with distinct
    gaps (e.g., number of children) {[}108{]}.
  \item
    \emph{Dependent vs.~Independent}: If variable \(Y\) is a consequence
    of variable \(X\), \(Y\) is dependent and \(X\) is independent
    {[}108{]}.
  \end{itemize}
\item
  \textbf{Extraneous Variable}: Independent variables not related to the
  study's purpose but affecting the dependent variable (causing
  experimental error) {[}108{]}.
\item
  \textbf{Control}: Minimizing the influence of extraneous variables
  using design restraints {[}108{]}.
\item
  \textbf{Confounded Relationship}: When the dependent variable's change
  is not free from the influence of extraneous variables {[}108{]}.
\item
  \textbf{Research Hypothesis}: A predictive statement that relates an
  independent variable to a dependent variable to be tested {[}109{]}.
\item
  \textbf{Experimental and Non-experimental Hypothesis-Testing
  Research}:

  \begin{itemize}
  \tightlist
  \item
    \emph{Experimental}: The independent variable is manipulated
    {[}109{]}.
  \item
    \emph{Non-experimental}: The independent variable is NOT manipulated
    (e.g., correlation of IQ and reading ability) {[}109{]}.
  \end{itemize}
\item
  \textbf{Experimental \& Control Groups}: Exposed to special/novel
  conditions (experimental) vs.~usual/baseline conditions (control)
  {[}110{]}.
\item
  \textbf{Treatments}: The different conditions under which groups are
  placed {[}110{]}.
\item
  \textbf{Experiment}: Process of examining the truth of a statistical
  hypothesis (absolute or comparative) {[}110{]}.
\item
  \textbf{Experimental Unit(s)}: Pre-determined plots or blocks where
  treatments are applied {[}111{]}.
\end{itemize}

\subsubsection{Different Research
Designs}\label{different-research-designs}

\begin{itemize}
\tightlist
\item
  \textbf{Exploratory (Formulative) Studies}: Focuses on discovering
  ideas and insights. Design must be flexible {[}112{]}. Three main
  methods:

  \begin{enumerate}
  \def\labelenumi{\arabic{enumi}.}
  \tightlist
  \item
    \emph{Survey of concerning literature}: Reviewing existing
    theoretical/empirical works {[}113{]}.
  \item
    \emph{Experience survey}: Interviewing competent people with
    practical experience {[}114{]}.
  \item
    \emph{Analysis of `insight-stimulating' examples}: Intensive study
    of selected striking instances {[}115{]}.
  \end{enumerate}
\item
  \textbf{Descriptive and Diagnostic Studies}:

  \begin{itemize}
  \tightlist
  \item
    \emph{Descriptive}: Portrays characteristics of individuals, groups,
    or situations {[}117{]}.
  \item
    \emph{Diagnostic}: Determines the frequency of something or its
    association with something else {[}117{]}.
  \item
    \emph{Design}: Must be rigid (not flexible) to protect against bias
    and maximize reliability. Relies on structured instruments and
    pre-planned statistical computations {[}117, 119, 120{]}.
  \end{itemize}
\item
  \textbf{Hypothesis-Testing Studies}: Experimental designs permitting
  causal inferences while minimizing bias {[}106{]}.
\end{itemize}

\subsubsection{Basic Principles of Experimental
Designs}\label{basic-principles-of-experimental-designs}

\begin{itemize}
\tightlist
\item
  \textbf{Fisher's Three Principles}:

  \begin{enumerate}
  \def\labelenumi{\arabic{enumi}.}
  \tightlist
  \item
    \emph{Principle of Replication}: Repeating the experiment more than
    once to apply treatments to multiple units, thereby increasing
    statistical accuracy and precision {[}123{]}.
  \item
    \emph{Principle of Randomization}: Randomly assigning treatments to
    experimental units to safeguard against extraneous variables and
    assure objective probability measurements {[}124{]}.
  \item
    \emph{Principle of Local Control}: Dividing units into homogeneous
    blocks and varying extraneous factors deliberately to measure and
    eliminate soil/block variability from experimental error {[}124{]}.
  \end{enumerate}
\end{itemize}

\subsubsection{Conclusion}\label{conclusion-1}

\begin{itemize}
\tightlist
\item
  The selection of design must be decided in advance based on the
  universe type, study objective, sampling frame, and desired accuracy
  {[}137{]}.
\end{itemize}

\subsubsection{Appendix: Developing a Research
Plan}\label{appendix-developing-a-research-plan}

\begin{itemize}
\tightlist
\item
  \textbf{Research Plan Items}: Clear statement of objectives, explicit
  description of the problem, operational definitions of concepts,
  method/approach to be used, detailed technical procedures
  (quantification details), population to be studied (sampling frame),
  methods of data processing/analysis, results of pilot tests, and
  time/cost budgets {[}141, 142, 143, 144{]}.
\end{itemize}

\begin{center}\rule{0.5\linewidth}{0.5pt}\end{center}

\subsection{Chapter 4: Sampling Design}\label{chapter-4-sampling-design}

\subsubsection{Census and Sample Survey}\label{census-and-sample-survey}

\begin{itemize}
\tightlist
\item
  \textbf{Census}: Complete enumeration of all items in the population.
  It is expensive, time-consuming, prone to non-sampling bias, and often
  physically impossible {[}145{]}.
\item
  \textbf{Sample Survey}: Examining a representative subset (sample) of
  the population. Yields sufficiently accurate results with less cost,
  time, and energy {[}145{]}.
\end{itemize}

\subsubsection{Implications of a Sample
Design}\label{implications-of-a-sample-design}

\begin{itemize}
\tightlist
\item
  \textbf{Sampling Design}: A definite plan determined before data
  collection to obtain a sample from a given population {[}127{]}.
\end{itemize}

\subsubsection{Steps in Sampling Design}\label{steps-in-sampling-design}

\begin{enumerate}
\def\labelenumi{\arabic{enumi}.}
\tightlist
\item
  \emph{Type of Universe}: Define the set of objects to be studied
  (Finite vs.~Infinite) {[}146{]}.
\item
  \emph{Sampling Unit}: Choose geographic (district), social (family),
  construction (house), or individual units {[}147{]}.
\item
  \emph{Source List (Sampling Frame)}: Prepare/acquire a reliable,
  comprehensive, and representative list of items of the universe
  {[}147{]}.
\item
  \emph{Size of Sample}: Choose an optimum sample size ensuring
  efficiency, representativeness, reliability, and cost-effectiveness
  {[}148{]}.
\item
  \emph{Parameters of Interest}: Determine whether estimating population
  mean, variance, or proportion {[}148{]}.
\item
  \emph{Budgetary Constraint}: Account for funds available for sampling
  {[}148{]}.
\item
  \emph{Sampling Procedure}: Decide the actual type of sample selection
  {[}127{]}.
\end{enumerate}

\subsubsection{Criteria of Selecting a Sampling
Procedure}\label{criteria-of-selecting-a-sampling-procedure}

\begin{itemize}
\tightlist
\item
  \textbf{Errors in Sampling}: Minimize systematic bias (caused by
  inappropriate frame, non-response, measurement errors) and sampling
  errors (random fluctuations) {[}128, 141{]}.
\end{itemize}

\subsubsection{Characteristics of a Good Sample
Design}\label{characteristics-of-a-good-sample-design}

\begin{itemize}
\tightlist
\item
  Must result in a truly representative sample, have small sampling
  error, be economically viable, control systematic bias, and allow
  generalization {[}128{]}.
\end{itemize}

\subsubsection{Different Types of Sample
Designs}\label{different-types-of-sample-designs}

\begin{itemize}
\tightlist
\item
  \textbf{Probability vs.~Non-probability Sampling}:

  \begin{itemize}
  \tightlist
  \item
    \emph{Non-probability}: Deliberate selection based on convenience or
    judgement. Probaiblity of selection is unknown (Haphazard,
    convenience, purposive, quota, judgement sampling) {[}59{]}.
  \item
    \emph{Probability (Random/Chance)}: Every item has an equal and
    known chance of selection (Simple random sampling, complex random
    sampling) {[}149{]}.
  \end{itemize}
\end{itemize}

\subsubsection{How to Select a Random
Sample?}\label{how-to-select-a-random-sample}

\begin{itemize}
\tightlist
\item
  \textbf{Lottery Method}: Draw folded slips from a container {[}149{]}.
\item
  \textbf{Random Number Tables}: Use Tippett's or other standardized
  tables to read numbers sequentially {[}150{]}.
\end{itemize}

\subsubsection{Random Sample from an Infinite
Universe}\label{random-sample-from-an-infinite-universe}

\begin{itemize}
\tightlist
\item
  Drawn by maintaining control over selection processes where each drawn
  item does not affect the infinite pool's properties {[}11{]}.
\end{itemize}

\subsubsection{Complex Random Sampling
Designs}\label{complex-random-sampling-designs}

\begin{itemize}
\tightlist
\item
  \textbf{Systematic Sampling}: Selecting every \(k^{th}\) item after a
  random start {[}130{]}.
\item
  \textbf{Stratified Sampling}: Dividing a heterogeneous population into
  homogeneous, non-overlapping subpopulations (strata) and drawing a
  random sample from each stratum {[}155{]}.

  \begin{itemize}
  \tightlist
  \item
    \emph{Proportional}: Strata sample size proportional to strata
    population size {[}132{]}.
  \item
    \emph{Optimum (Disproportionate)}: Strata sample size based on
    standard deviation and cost (\(n_i \propto N_i \sigma_i\))
    {[}152{]}.
  \end{itemize}
\item
  \textbf{Cluster \& Area Sampling}: Dividing population into clusters,
  randomly selecting clusters, and completely enumerating all units
  within chosen clusters {[}137{]}.
\item
  \textbf{Multi-stage Sampling}: Sampling in successive stages (e.g.,
  district \(ightarrow\) town \(ightarrow\) household) {[}138{]}.
\item
  \textbf{Sequential Sampling}: The sample size is not fixed in advance;
  sampling continues based on cumulative intermediate results {[}140{]}.
\end{itemize}

\begin{center}\rule{0.5\linewidth}{0.5pt}\end{center}

\subsection{Chapter 5: Measurement and Scaling
Techniques}\label{chapter-5-measurement-and-scaling-techniques}

\subsubsection{Measurement in Research}\label{measurement-in-research}

\begin{itemize}
\tightlist
\item
  \textbf{Definition}: The process of assigning numbers to physical
  objects or abstract concepts according to rules of correspondence
  mapping a domain onto a range {[}155, 156{]}.
\end{itemize}

\subsubsection{Measurement Scales}\label{measurement-scales}

\begin{itemize}
\tightlist
\item
  \textbf{Four Types}:

  \begin{enumerate}
  \def\labelenumi{\arabic{enumi}.}
  \tightlist
  \item
    \emph{Nominal}: Categorical classification with no quantitative
    value or order (e.g., Gender: 0 for male, 1 for female) {[}156{]}.
  \item
    \emph{Ordinal}: Ranks items in order of magnitude (relative
    positions) but distances between positions are unequal or unknown
    (e.g., military ranks) {[}152{]}.
  \item
    \emph{Interval}: Ranks with equal, standardized intervals between
    positions but has an arbitrary zero point (e.g., Fahrenheit
    temperature) {[}152{]}.
  \item
    \emph{Ratio}: Ranks with equal intervals and an absolute, true zero
    point representing absence of the trait (e.g., weight, height)
    {[}152{]}.
  \end{enumerate}
\end{itemize}

\subsubsection{Sources of Error in
Measurement}\label{sources-of-error-in-measurement}

\begin{itemize}
\tightlist
\item
  \textbf{Four Factors}:

  \begin{enumerate}
  \def\labelenumi{\arabic{enumi}.}
  \tightlist
  \item
    \emph{Respondent}: Fatigue, boredom, anxiety, or ignorance
    {[}161{]}.
  \item
    \emph{Situation}: Presence of others, noise, or lack of privacy
    {[}161{]}.
  \item
    \emph{Measurer}: Interviewer bias, rewording, or improper coding
    {[}120, 161{]}.
  \item
    \emph{Instrument}: Defective layout, ambiguous wording, or
    formatting issues {[}162{]}.
  \end{enumerate}
\end{itemize}

\subsubsection{Tests of Sound
Measurement}\label{tests-of-sound-measurement}

\begin{itemize}
\tightlist
\item
  \textbf{Three Key Criteria}:

  \begin{enumerate}
  \def\labelenumi{\arabic{enumi}.}
  \tightlist
  \item
    \emph{Validity}: Extent to which an instrument measures what it
    claims to measure {[}157{]}.

    \begin{itemize}
    \tightlist
    \item
      \emph{Content}: Adequate coverage of the topic
      (judgmental/intuitive) {[}158{]}.
    \item
      \emph{Criterion-Related (Predictive/Concurrent)}: Ability to
      predict or relate to an external criterion {[}158, 159{]}.
    \item
      \emph{Construct}: Extent to which it measures a theoretical trait
      {[}158{]}.
    \end{itemize}
  \item
    \emph{Reliability}: Stability (consistent over time) and Equivalence
    (consistent across items/investigators) {[}160{]}.
  \item
    \emph{Practicality}: Judged by economy (budget trade-offs),
    convenience (easy layout), and interpretability (supplemented by
    manual and scoring keys) {[}162{]}.
  \end{enumerate}
\end{itemize}

\subsubsection{Technique of Developing Measurement
Tools}\label{technique-of-developing-measurement-tools}

\begin{itemize}
\tightlist
\item
  \textbf{Four Stages}: (a) Concept development, (b) Specification of
  concept dimensions, (c) Selection of indicators, and (d) Formation of
  an index {[}163{]}.
\end{itemize}

\subsubsection{Scaling}\label{scaling}

\begin{itemize}
\tightlist
\item
  \textbf{Definition}: Procedures of assigning numbers to degrees of
  opinions, attitudes, or institutional concepts along a continuum
  {[}167{]}.
\end{itemize}

\subsubsection{Scale Classification
Bases}\label{scale-classification-bases}

\begin{itemize}
\tightlist
\item
  Classified by: (a) Subject orientation, (b) Response form, (c) Degree
  of subjectivity, (d) Scale properties, (e) Number of dimensions, and
  (f) Scale construction techniques {[}168{]}.
\end{itemize}

\subsubsection{Important Scaling
Techniques}\label{important-scaling-techniques}

\begin{itemize}
\tightlist
\item
  \textbf{Rating Scales}: Absolute evaluation against categories without
  reference to other objects (e.g., graphic scale or itemized scale)
  {[}171{]}.
\item
  \textbf{Ranking Scales}: Relative evaluation comparing objects
  directly (e.g., Paired Comparison, Method of Equal-Appearing
  Intervals, or Rank Order) {[}171, 172{]}.
\end{itemize}

\subsubsection{Scale Construction
Techniques}\label{scale-construction-techniques}

\begin{itemize}
\tightlist
\item
  \textbf{Five Major Approaches}:

  \begin{enumerate}
  \def\labelenumi{\arabic{enumi}.}
  \tightlist
  \item
    \emph{Arbitrary}: Ad-hoc, subjective item selection {[}169{]}.
  \item
    \emph{Consensus (Thurstone)}: Submitted to judges who sort
    statements into 11 piles. Statements with high agreement (low
    variance) and evenly spread median values constitute the final scale
    {[}177, 178{]}.
  \item
    \emph{Item Analysis (Likert Summated)}: 5-point agreements scored.
    Extreme top and bottom 25\% groups act as criteria; items with high
    discriminatory power are retained {[}179, 180, 181, 182{]}.
  \item
    \emph{Cumulative (Guttman Scalogram)}: Unidimensional, structured
    series where agreement with an extreme item implies agreement with
    all less extreme items. Evaluated by Coefficient of Reproducibility
    (\(C_R = 1 - e/(nN)\); minimum standard 0.90) {[}185, 186, 191{]}.
  \item
    \emph{Factor Scales}: Based on multidimensional intercorrelations
    {[}192{]}.

    \begin{itemize}
    \tightlist
    \item
      \emph{Semantic Differential}: Charles Osgood's 7-point bipolar
      ratings evaluating three dominant dimensions: Evaluation, Potency,
      and Activity {[}193, 194{]}.
    \item
      \emph{Multidimensional Scaling (MDS)}: Plots judged similarities
      as physical distances in a multidimensional geometric space
      {[}198{]}.
    \end{itemize}
  \end{enumerate}
\end{itemize}

\begin{center}\rule{0.5\linewidth}{0.5pt}\end{center}

\subsection{Chapter 6: Methods of Data
Collection}\label{chapter-6-methods-of-data-collection}

\subsubsection{Collection of Primary
Data}\label{collection-of-primary-data}

\begin{itemize}
\tightlist
\item
  \textbf{Primary Data}: Collected afresh and for the first time;
  original in character {[}203{]}.
\end{itemize}

\subsubsection{Observation Method}\label{observation-method}

\begin{itemize}
\tightlist
\item
  \textbf{Definition}: Systematic viewing of raw behavior {[}96{]}.
\item
  \textbf{Types}: Structured vs.~Unstructured {[}96{]}, Participant
  vs.~Non-participant {[}96{]}, and Controlled vs.~Uncontrolled
  {[}97{]}.
\item
  \textbf{Evaluation}: Eliminates subjective bias and record independent
  of respondent willingness, but is expensive and limited to current
  behavior {[}96, 97{]}.
\end{itemize}

\subsubsection{Interview Method}\label{interview-method}

\begin{itemize}
\tightlist
\item
  \textbf{Personal Interview}: Structured (standardized set of
  questions) or Unstructured (flexible sequence/wording) {[}204{]}.
\item
  \textbf{Other Types}: Focussed (focuses on concrete experiences),
  Clinical (broad motivations), Non-directive (interviewer acts as a
  passive catalyst), and Telephone {[}204, 205{]}.
\item
  \textbf{Evaluation}: High response rate, deeper probing, but expensive
  and prone to interviewer bias {[}203{]}.
\end{itemize}

\subsubsection{Collection of Data through
Questionnaires}\label{collection-of-data-through-questionnaires}

\begin{itemize}
\tightlist
\item
  \textbf{Definition}: Mailed or sent series of typed questions to be
  completed independently by respondents {[}208{]}.
\item
  \textbf{Main Aspects}: General form (structured vs.~unstructured)
  {[}209{]}, Question sequence (general-to-specific; vital questions in
  middle) {[}210, 211{]}, and Question formulation (simple, concrete,
  unambiguous, avoiding ``danger/catch'' words) {[}211, 213{]}.
\item
  \textbf{Wording Forms}: Multiple choice (closed) vs.~Open-ended
  (latitude in phrasing) {[}212{]}.
\item
  \textbf{Essentials}: Short, attractive physical appearance, clear
  instructions, and includes control questions for verification {[}214,
  215{]}.
\end{itemize}

\subsubsection{Collection of Data through
Schedules}\label{collection-of-data-through-schedules}

\begin{itemize}
\tightlist
\item
  \textbf{Definition}: Questionnaires filled out in person by trained
  enumerators {[}104{]}. Enumerators must be honest, patient, and
  trained to cross-examine {[}216{]}.
\end{itemize}

\subsubsection{Difference between Questionnaires and
Schedules}\label{difference-between-questionnaires-and-schedules}

\begin{itemize}
\tightlist
\item
  \emph{Questionnaires}: Cheaper, sent by mail, relies on respondent
  literacy, suffers from high non-response, and has no direct contact
  {[}214{]}.
\item
  \emph{Schedules}: Costly, filled out by enumerators, can use
  illiterate respondents, low non-response, and permits identity
  verification {[}214, 216{]}.
\end{itemize}

\subsubsection{Some Other Methods of Data
Collection}\label{some-other-methods-of-data-collection}

\begin{itemize}
\tightlist
\item
  \textbf{Pantry \& Store Audits}: Physical count of stocks {[}106{]}.
\item
  \textbf{Consumer Panels}: Transcribe continuous diaries of purchases
  {[}106{]}.
\item
  \textbf{Projective Techniques}: Indirect assessment of unconscious
  motivations {[}108{]}.

  \begin{itemize}
  \tightlist
  \item
    \emph{Word Association \& Sentence/Story Completion} {[}217, 218{]}.
  \item
    \emph{TAT \& HIT (Holtzman Inkblot Test)}: HIT uses 45 cards with
    single responses to assess barrier scores {[}219, 221{]}.
  \item
    \emph{Play \& Sociometry (sociological patterns)} {[}222{]}.
  \end{itemize}
\item
  \textbf{Depth Interviews \& Content Analysis}: Motivational probing
  vs.~systematic qualitative evaluation of documents {[}223, 224{]}.
\end{itemize}

\subsubsection{Collection of Secondary
Data}\label{collection-of-secondary-data}

\begin{itemize}
\tightlist
\item
  \textbf{Secondary Data}: Already collected and analyzed by someone
  else {[}226{]}.
\item
  \textbf{Criteria for Use}:

  \begin{enumerate}
  \def\labelenumi{\arabic{enumi}.}
  \tightlist
  \item
    \emph{Reliability}: Who collected it? What methods? Was there bias?
    {[}227{]}
  \item
    \emph{Suitability}: Do definitions and scope match the present
    study? {[}227{]}
  \item
    \emph{Adequacy}: Accuracy levels and physical coverage boundaries
    {[}228{]}.
  \end{enumerate}
\end{itemize}

\subsubsection{Selection of Appropriate Method for Data
Collection}\label{selection-of-appropriate-method-for-data-collection}

\begin{itemize}
\tightlist
\item
  Depends on the: (i) Nature, scope, and object of enquiry, (ii)
  Availability of funds, (iii) Time factor, and (iv) Precision required
  {[}229{]}.
\end{itemize}

\subsubsection{Case Study Method}\label{case-study-method}

\begin{itemize}
\tightlist
\item
  \textbf{Definition}: Intensive qualitative observation and in-depth
  analysis of a single social unit (person, family, institution,
  community) as an integrated totality {[}230, 231{]}.
\item
  \textbf{Phases}: Locating the status, collecting history, diagnosing
  causal factors, applying therapy, and follow-up {[}234, 235{]}.
\item
  \textbf{Assumptions}: Uniformity in basic human nature; completeness
  of historical sequence {[}234{]}.
\end{itemize}

\begin{center}\rule{0.5\linewidth}{0.5pt}\end{center}

\subsection{Chapter 7: Processing and Analysis of
Data}\label{chapter-7-processing-and-analysis-of-data}

\subsubsection{Processing Operations}\label{processing-operations}

\begin{itemize}
\tightlist
\item
  \textbf{Editing}: Scrutinizing completed schedules/questionnaires to
  detect and correct errors and omissions.

  \begin{itemize}
  \tightlist
  \item
    \emph{Field Editing}: Completed by investigator immediately on
    same/next day {[}252{]}.
  \item
    \emph{Central Editing}: Thorough, centralized check by a single
    editor/team {[}253{]}.
  \end{itemize}
\item
  \textbf{Coding}: Assigning numbers/symbols to categories of answers.
  Categories must be exhaustive, mutually exclusive, and unidimensional
  {[}255, 256{]}.
\item
  \textbf{Classification}: Dividing data into homogeneous groups.

  \begin{itemize}
  \tightlist
  \item
    \emph{By Attributes (Descriptive)}: Qualitative presence/absence
    (simple or manifold classification) {[}258{]}.
  \item
    \emph{By Class Intervals (Numerical)}: Quantitative measures. Size
    of interval (\(i\)) determined using Sturges' Formula
    (\(i = R/(1 + 3.3 \log N)\)) {[}259, 262{]}. Limits can be
    \emph{exclusive} (upper limit excluded) {[}264{]} or
    \emph{inclusive} (upper limit included) {[}265{]}.
  \end{itemize}
\item
  \textbf{Tabulation}: Orderly arrangement of summarized data in columns
  and rows {[}267{]}. Requires a distinct table number, clear title,
  captions (column headings), stubs (row headings), units, and footnotes
  {[}269, 270{]}.
\end{itemize}

\subsubsection{Some Problems in
Processing}\label{some-problems-in-processing}

\begin{itemize}
\tightlist
\item
  \textbf{Don't Know (DK) Responses}: Solved by design adjustments,
  allocation from other data, or proportionate distribution {[}274,
  275{]}.
\item
  \textbf{Percentages}: Reducing data to a 0-100 base. Never average
  separate percentages unless weighted by group size; higher figures are
  base for decreases {[}275, 276{]}.
\end{itemize}

\subsubsection{Elements/Types of
Analysis}\label{elementstypes-of-analysis}

\begin{itemize}
\tightlist
\item
  \textbf{Descriptive}: Largely the study of distributions of one
  variable (unidimensional), two variables (bivariate), or more
  (multivariate) {[}277{]}.
\item
  \textbf{Inferential (Statistical)}: Testing hypotheses to determine
  validity and make generalizations {[}279{]}.
\end{itemize}

\subsubsection{Statistics in Research}\label{statistics-in-research}

\begin{itemize}
\tightlist
\item
  \textbf{Descriptive Statistics}: Development of indices from raw data
  {[}279{]}.
\item
  \textbf{Inferential Statistics}: Generalization from sample to
  population parameters {[}279{]}.
\end{itemize}

\subsubsection{Measures of Central
Tendency}\label{measures-of-central-tendency}

\begin{itemize}
\tightlist
\item
  \textbf{Arithmetic Mean}: Sum of values divided by \(n\) (simple or
  weighted; stable, algebraic, but affected by extreme items) {[}280,
  281{]}.
\item
  \textbf{Median}: Positional middle value of an arrayed series
  (unaffected by extremes) {[}285{]}.
\item
  \textbf{Mode}: Most frequently occurring value; concentration peak
  (useful for popular sizes) {[}282{]}.
\item
  \textbf{Geometric \& Harmonic Means}: Reciprocals of arithmetic
  averages of logarithms and reciprocals {[}133, 283{]}.
\end{itemize}

\subsubsection{Measures of Dispersion}\label{measures-of-dispersion}

\begin{itemize}
\tightlist
\item
  \textbf{Range}: Difference between largest and smallest item
  {[}134{]}.
\item
  \textbf{Mean Deviation}: Average of absolute deviations from
  mean/median/mode {[}284{]}.
\item
  \textbf{Standard Deviation (\(\sigma\))}: Root-mean-square deviation.
  Gives the measure of dispersion and basis for standard error {[}135,
  303{]}.
\item
  \textbf{Coefficient of Variation}: Ratio of standard deviation to mean
  multiplied by 100 {[}136{]}.
\end{itemize}

\subsubsection{Measures of Asymmetry
(Skewness)}\label{measures-of-asymmetry-skewness}

\begin{itemize}
\tightlist
\item
  Indicates shape of distribution and divergence from normal symmetry
  {[}136{]}.
\end{itemize}

\subsubsection{Measures of Relationship}\label{measures-of-relationship}

\begin{itemize}
\tightlist
\item
  \textbf{Simple Correlation (Karl Pearson's \(r\))}: Linear
  relationship magnitude (-1 to +1) {[}139, 298{]}.
\item
  \textbf{Partial \& Multiple Correlation}: Relationship of variables
  when others are controlled or combined {[}142, 143{]}.
\end{itemize}

\subsubsection{Simple Regression
Analysis}\label{simple-regression-analysis}

\begin{itemize}
\tightlist
\item
  Formulating a mathematical model (\(Y = a + bX\)) for predicting a
  dependent variable based on independent values {[}285{]}.
\end{itemize}

\subsubsection{Multiple Correlation and
Regression}\label{multiple-correlation-and-regression}

\begin{itemize}
\tightlist
\item
  Simultaneous assessment of multiple independent variables on a single
  metric dependent variable {[}142, 143{]}.
\end{itemize}

\subsubsection{Association in Case of
Attributes}\label{association-in-case-of-attributes}

\begin{itemize}
\tightlist
\item
  Tested via Yule's Coefficient of Association (\(Q\)) {[}286{]}.
  Illusory association occurs when real relationship is absent
  {[}287{]}. Tested on \(r \times c\) contingency tables {[}288{]}.
\end{itemize}

\subsubsection{Other Measures}\label{other-measures}

\begin{itemize}
\tightlist
\item
  \textbf{Index Numbers}: Approximate indicators of directional trend
  change over time {[}290{]}.
\item
  \textbf{Time Series Analysis}: Decomposing data into Trend (\(T\)),
  Cyclical (\(C\)), Seasonal (\(S\)), and Irregular (\(I\)) components
  {[}291{]}.
\end{itemize}

\begin{center}\rule{0.5\linewidth}{0.5pt}\end{center}

\subsection{Chapter 8: Sampling
Fundamentals}\label{chapter-8-sampling-fundamentals}

\subsubsection{Need for Sampling}\label{need-for-sampling}

\begin{itemize}
\tightlist
\item
  Enables rapid results, reduces costs, controls non-sampling errors,
  and permits destructive testing {[}152, 300{]}.
\end{itemize}

\subsubsection{Some Fundamental
Definitions}\label{some-fundamental-definitions}

\begin{itemize}
\tightlist
\item
  \textbf{Universe/Population}: Total items in field of inquiry
  (elementary units) {[}300{]}.
\item
  \textbf{Statistic vs.~Parameter}: Sample-derived indices
  vs.~Population-derived values {[}311{]}.
\item
  \textbf{Sampling Frame}: Source list of population items {[}311{]}.
\item
  \textbf{Sampling Error}: Fluctuations due to random chance {[}311{]}.
\end{itemize}

\subsubsection{Important Sampling
Distributions}\label{important-sampling-distributions}

\begin{itemize}
\tightlist
\item
  Normal, Student's \(t\), Chi-square, and \(F\) distributions,
  depending on parameters and sample sizes {[}155, 311{]}.
\end{itemize}

\subsubsection{Central Limit Theorem}\label{central-limit-theorem}

\begin{itemize}
\tightlist
\item
  Sample means of samples drawn from \emph{any} population approach
  normal distribution as sample size \(n\) increases (\(n > 30\))
  {[}301{]}.
\end{itemize}

\subsubsection{Sampling Theory}\label{sampling-theory}

\begin{itemize}
\tightlist
\item
  Deals with math of sampling accuracy and standard error based on laws
  of statistical regularity and inertia of large numbers {[}158{]}.
\end{itemize}

\subsubsection{Sandler's A-test}\label{sandlers-a-test}

\begin{itemize}
\tightlist
\item
  A paired-sample shortcut alternative to Student's \(t\)-test based on
  difference scores: \(A = \sum D_i^2 / (\sum D_i)^2\) {[}302{]}.
\end{itemize}

\subsubsection{Concept of Standard
Error}\label{concept-of-standard-error}

\begin{itemize}
\tightlist
\item
  The standard error is the standard deviation of the sampling
  distribution. Smaller S.E. indicates greater reliability {[}303{]}.
\end{itemize}

\subsubsection{Estimation}\label{estimation}

\begin{itemize}
\tightlist
\item
  \textbf{Point Estimate}: A single value estimate {[}311{]}.
\item
  \textbf{Interval Estimate}: A range of values within confidence limits
  {[}311{]}.
\end{itemize}

\subsubsection{\texorpdfstring{Estimating the Population Mean
(\(\mu\))}{Estimating the Population Mean (\textbackslash mu)}}\label{estimating-the-population-mean-mu}

\begin{itemize}
\tightlist
\item
  Calculated as \(bar{X} \pm z \cdot S.E.\) using finite population
  multipliers if needed {[}305, 336{]}.
\end{itemize}

\subsubsection{Estimating Population
Proportion}\label{estimating-population-proportion}

\begin{itemize}
\tightlist
\item
  Calculated as \(p \pm z \cdot \sqrt{pq/n}\) for attribute frequencies
  {[}172{]}.
\end{itemize}

\subsubsection{Sample Size and its
Determination}\label{sample-size-and-its-determination}

\begin{itemize}
\tightlist
\item
  Depends on: Universe nature (heterogeneity), number of proposed
  classes, precision, and confidence level {[}305, 306{]}.
\end{itemize}

\subsubsection{Determination of Sample Size: Precision \& Confidence
Approach}\label{determination-of-sample-size-precision-confidence-approach}

\begin{itemize}
\tightlist
\item
  \textbf{Means}: \(n = z^2 \sigma_p^2 / e^2\) {[}307{]}.
\item
  \textbf{Proportions}: \(n = z^2 p q / e^2\) {[}308{]}.
\end{itemize}

\subsubsection{Determination of Sample Size: Bayesian
Approach}\label{determination-of-sample-size-bayesian-approach}

\begin{itemize}
\tightlist
\item
  Maximizes Expected Net Gain (ENG) by balancing Expected Value of
  Sample Information (EVSI) against cost of sampling {[}309{]}.
\end{itemize}

\begin{center}\rule{0.5\linewidth}{0.5pt}\end{center}

\subsection{Chapter 9: Testing of Hypotheses-I (Parametric
Tests)}\label{chapter-9-testing-of-hypotheses-i-parametric-tests}

\subsubsection{What is a Hypothesis?}\label{what-is-a-hypothesis}

\begin{itemize}
\tightlist
\item
  \textbf{Definition}: A tentative predictive statement or assumption
  made to test logical or empirical consequences {[}317{]}.
\item
  \textbf{Characteristics}: Clear, testable, specific, simple, limited
  in scope, and consistent with known facts {[}318, 319, 320{]}.
\end{itemize}

\subsubsection{Basic Concepts Concerning Testing of
Hypotheses}\label{basic-concepts-concerning-testing-of-hypotheses}

\begin{itemize}
\tightlist
\item
  \textbf{Null (\(H_0\)) vs.~Alternative (\(H_a\))}: Hypothesis of no
  difference vs.~the operational alternative to prove {[}321, 323{]}.
\item
  \textbf{Level of Significance (\(\alpha\))}: Maximum probability of
  Type I error (rejecting true \(H_0\)) chosen in advance {[}324,
  325{]}.
\item
  \textbf{Type I vs.~Type II Errors}:

  \begin{itemize}
  \tightlist
  \item
    \emph{Type I (\(\alpha\))}: Rejecting true \(H_0\) {[}327{]}.
  \item
    \emph{Type II (\(\beta\))}: Accepting false \(H_0\) {[}327{]}.
  \end{itemize}
\item
  \textbf{Two-tailed vs.~One-tailed Tests}: Rejection region on both
  tails (\(H_a: \mu \neq \mu_0\)) vs.~single tail (\(H_a: \mu > \mu_0\))
  {[}328{]}.
\end{itemize}

\subsubsection{Procedure for Hypothesis
Testing}\label{procedure-for-hypothesis-testing}

\begin{itemize}
\tightlist
\item
  \textbf{Sequence}: Formulate \(H_0/H_a \rightarrow\) specify
  \(\alpha ightarrow\) select sampling distribution \(ightarrow\)
  calculate test statistic \(ightarrow\) find probability \(ightarrow\)
  accept or reject \(H_0\) {[}331{]}.
\end{itemize}

\subsubsection{Important Parametric
Tests}\label{important-parametric-tests}

\begin{itemize}
\tightlist
\item
  \textbf{z-test}: \(n > 30\) or known population variance {[}333,
  334{]}.
\item
  \textbf{t-test}: \(n \le 30\) and unknown population variance
  {[}335{]}.
\item
  \textbf{F-test}: Compares the variance of two normal populations
  {[}333{]}.
\item
  \textbf{Sandler's A-test}: Shortcut pairing test {[}335{]}.
\end{itemize}

\subsubsection{Limitations of the Tests of
Hypotheses}\label{limitations-of-the-tests-of-hypotheses}

\begin{itemize}
\tightlist
\item
  Tests are not absolute proofs; they are tools of decision-making
  subject to probability limitations and require judgmental context
  {[}358{]}.
\end{itemize}

\begin{center}\rule{0.5\linewidth}{0.5pt}\end{center}

\subsection{Chapter 10: Chi-Square
Test}\label{chapter-10-chi-square-test}

\subsubsection{Chi-Square as a Test for Comparing
Variance}\label{chi-square-as-a-test-for-comparing-variance}

\begin{itemize}
\tightlist
\item
  Tests whether sample variance matches a hypothesized population
  variance: \(\chi^2 = \sigma_s^2 (n - 1) / \sigma_p^2\) {[}16, 367{]}.
\end{itemize}

\subsubsection{Chi-Square as a Non-parametric
Test}\label{chi-square-as-a-non-parametric-test}

\begin{itemize}
\tightlist
\item
  \textbf{Goodness of Fit}: Compares observed (\(O\)) and expected
  (\(E\)) frequencies: \(\chi^2 = \sum (O_i - E_i)^2 / E_i\) {[}301{]}.
\item
  \textbf{Independence}: Evaluates whether two attributes in a
  contingency table are associated {[}333{]}.
\end{itemize}

\subsubsection{\texorpdfstring{Conditions for the Application of
\(\chi^2\)
Test}{Conditions for the Application of \textbackslash chi\^{}2 Test}}\label{conditions-for-the-application-of-chi2-test}

\begin{enumerate}
\def\labelenumi{\arabic{enumi}.}
\tightlist
\item
  Observations randomly collected {[}368{]}.
\item
  Items in sample independent {[}368{]}.
\item
  Expected frequencies in any cell must be \(\ge 10\) (combine adjacent
  cells if smaller; some accept \(\ge 5\)) {[}368{]}.
\item
  Total \(N \ge 50\) {[}369{]}.
\end{enumerate}

\subsubsection{Yates' Correction}\label{yates-correction}

\begin{itemize}
\tightlist
\item
  Applied to \(2 \times 2\) tables with expected cell frequencies
  \(< 5\) by subtracting \(0.5\) from absolute differences:
  \(\chi^2_{adj} = \sum (|O - E| - 0.5)^2 / E\) {[}16, 655{]}.
\end{itemize}

\begin{center}\rule{0.5\linewidth}{0.5pt}\end{center}

\subsection{Chapter 11: Analysis of Variance and
Covariance}\label{chapter-11-analysis-of-variance-and-covariance}

\subsubsection{Analysis of Variance
(ANOVA)}\label{analysis-of-variance-anova}

\begin{itemize}
\tightlist
\item
  \textbf{Definition}: Decomposes total variance into component parts
  attributable to specific causes vs.~random error {[}361, 364{]}.
\end{itemize}

\subsubsection{The Basic Principle of
ANOVA}\label{the-basic-principle-of-anova}

\begin{itemize}
\tightlist
\item
  Compares variance between treatment groups against variance within
  groups using the \(F\)-ratio: \(F = MS_{between} / MS_{within}\)
  {[}17, 381{]}.
\end{itemize}

\subsubsection{ANOVA Technique}\label{anova-technique}

\begin{itemize}
\tightlist
\item
  \textbf{One-way ANOVA}: Compares \(k\) sample means {[}382{]}.
\item
  \textbf{Two-way ANOVA}: Compares treatment effects and controls a
  metric covariate, evaluating interaction effects {[}127, 392{]}.
\item
  \textbf{ANOVA in Latin-Square Design}: Controls two extraneous sources
  of variation (e.g., soil fertility rows and seed columns) {[}129,
  398{]}.
\end{itemize}

\subsubsection{Analysis of Co-variance
(ANOCOVA)}\label{analysis-of-co-variance-anocova}

\begin{itemize}
\tightlist
\item
  Controls the influence of an uncontrolled metric independent variable
  (covariate) on the dependent variable using linear regression {[}391,
  402{]}.
\end{itemize}

\begin{center}\rule{0.5\linewidth}{0.5pt}\end{center}

\subsection{Chapter 12: Testing of Hypotheses-II (Nonparametric
Tests)}\label{chapter-12-testing-of-hypotheses-ii-nonparametric-tests}

\subsubsection{Important Nonparametric or Distribution-free
Tests}\label{important-nonparametric-or-distribution-free-tests}

\begin{itemize}
\tightlist
\item
  Do not rely on population normality assumptions {[}417{]}.
\item
  \textbf{Sign Test}: One-sample or two-sample matched pair evaluations
  based on signs of differences {[}418, 419{]}.
\item
  \textbf{McNemer Test}: Evaluates significance of change in nominal
  matching before-and-after views {[}423{]}.
\item
  \textbf{Wilcoxon Signed Rank Test}: Ranks absolute differences to
  check matched pairs {[}615{]}.
\item
  \textbf{Mann-Whitney U-test}: Compares medians of two independent
  groups using rank sums {[}665{]}.
\item
  \textbf{Kruskal-Wallis H-test}: Nonparametric one-way ANOVA utilizing
  rank sums {[}425, 427{]}.
\item
  \textbf{Runs Test}: Evaluates sequence randomness based on number of
  runs of binary attributes (\(r\)) {[}428{]}.
\item
  \textbf{Spearman Rank Correlation (\(r_s\))}: Measures ordinal
  association {[}430{]}.
\item
  \textbf{Kendall Concordance (\(W\))}: Measures consensus agreement
  among \(k\) judges ranking \(N\) items {[}433{]}.
\end{itemize}

\begin{center}\rule{0.5\linewidth}{0.5pt}\end{center}

\subsection{Chapter 13: Multivariate Analysis
Techniques}\label{chapter-13-multivariate-analysis-techniques}

\subsubsection{Classification of Multivariate
Techniques}\label{classification-of-multivariate-techniques}

\begin{itemize}
\tightlist
\item
  \textbf{Dependence}: Explanatory variables predict criterion
  variables.

  \begin{itemize}
  \tightlist
  \item
    \emph{Multiple Regression}: Metric criterion predicted from metric
    explanatory variables {[}450{]}.
  \item
    \emph{Discriminant Analysis}: Non-metric categorical group predicted
    from metric variables {[}452, 456{]}.
  \item
    \emph{MANOVA}: Multiple metric dependent variables evaluated against
    non-metric experimental factors {[}456{]}.
  \item
    \emph{Canonical Correlation}: Predicts a metric set from another
    metric set {[}456{]}.
  \end{itemize}
\item
  \textbf{Interdependence}: Evaluates mutual relationships without
  dependency.

  \begin{itemize}
  \tightlist
  \item
    \emph{Factor Analysis}: Groups correlated manifest variables into
    fewer latent factors {[}457, 458{]}.
  \item
    \emph{Cluster Analysis}: Classifies entities into homogeneous,
    distinct groups {[}496{]}.
  \item
    \emph{MDS}: Maps metric/non-metric perceptual similarities into
    geometric space {[}198, 500{]}.
  \end{itemize}
\item
  \textbf{Path Analysis}: Uses path diagrams to decompose total
  correlation into direct and indirect effects using standardized
  regression beta weights {[}501, 504{]}.
\end{itemize}

\begin{center}\rule{0.5\linewidth}{0.5pt}\end{center}

\subsection{Chapter 14: Interpretation and Report
Writing}\label{chapter-14-interpretation-and-report-writing}

\subsubsection{Meaning and Necessity of
Interpretation}\label{meaning-and-necessity-of-interpretation}

\begin{itemize}
\tightlist
\item
  Refers to drawing inferences from collected facts; links findings to
  existing theories and establishes continuity {[}507, 508{]}.
\end{itemize}

\subsubsection{Different Steps in Writing
Report}\label{different-steps-in-writing-report}

\begin{enumerate}
\def\labelenumi{\arabic{enumi}.}
\tightlist
\item
  Logical analysis of subject-matter (logical/chronological) {[}518{]}.
\item
  Preparation of final outline {[}518{]}.
\item
  Preparation of rough draft {[}519{]}.
\item
  Rewriting and polishing {[}519{]}.
\item
  Preparation of final bibliography (books and periodicals list) {[}519,
  520{]}.
\item
  Writing the final draft {[}522{]}.
\end{enumerate}

\subsubsection{Layout of the Research
Report}\label{layout-of-the-research-report}

\begin{itemize}
\tightlist
\item
  \textbf{Preliminary Pages}: Title, preface/foreword, Table of
  Contents, list of tables/charts {[}524{]}.
\item
  \textbf{Main Text}: Introduction, statement of
  findings/recommendations, results, implications, and summary
  {[}524{]}.
\item
  \textbf{End Matter}: Appendices, bibliography, and alphabetical
  subject/author index {[}530{]}.
\end{itemize}

\subsubsection{Types of Reports}\label{types-of-reports}

\begin{itemize}
\tightlist
\item
  \textbf{Technical Report}: Focuses on methods, assumptions, detailed
  data, and limitations {[}533, 534{]}.
\item
  \textbf{Popular Report}: Focuses on practical findings, simplicity,
  subheadings, and visual charts {[}537{]}.
\end{itemize}

\begin{center}\rule{0.5\linewidth}{0.5pt}\end{center}

\subsection{Chapter 15: The Computer: Its Role in
Research}\label{chapter-15-the-computer-its-role-in-research}

\subsubsection{The Computer System}\label{the-computer-system}

\begin{itemize}
\tightlist
\item
  \textbf{CPU Segments}: Internal Storage (holds programs/data), Control
  Unit (sequencing), and Arithmetic Logical Unit (ALU - calculations)
  {[}570{]}.
\item
  \textbf{Terminology}: Hardware, Software (programs), Firmware, and
  System/Application Software {[}573, 574{]}.
\end{itemize}

\subsubsection{Important
Characteristics}\label{important-characteristics}

\begin{itemize}
\tightlist
\item
  High speed, diligence (no fatigue), storage, accuracy, and automation
  {[}575, 576, 577{]}.
\end{itemize}

\subsubsection{The Binary Number System}\label{the-binary-number-system}

\begin{itemize}
\tightlist
\item
  Operates on base-2 (0 and 1). Arithmetic is reduced to rapid binary
  addition using complement subtraction {[}570, 581, 582{]}.
\end{itemize}

\subsubsection{Computers and Researcher}\label{computers-and-researcher}

\begin{itemize}
\tightlist
\item
  Enables rapid statistical analysis (Canned packages for regression,
  factor analysis, ANOVA) {[}587, 588{]}.
\item
  \emph{Limitation}: \textbf{GIGO} (``Garbage In, Garbage Out'') -- the
  quality of the analysis is entirely dependent on the quality of inputs
  and programming {[}592, 593{]}.
\end{itemize}

\end{document}
